\begin{lemma}[Incoming reconstruction from the third matching edge]
Fix an injective map $g:A\to B$, $x,a\in A$, and $y\in B$.
Suppose $f_1,f_2:A\to B$ and $b_1,b_2\in A$ satisfy
$f_1(x)=f_2(x)=y$, that $x,a,b_j$ are pairwise distinct for each
$j\in\{1,2\}$, and that
\[
 S(f_1;x,a,b_1)=g=S(f_2;x,a,b_2)
\]
with the switch of \noderef{MatchingThreeEdgeSwitch}.
Then $b_1=b_2$ and $f_1=f_2$. Thus, for a fixed outgoing matching $g$
and removed edge $xy$, the left endpoint $a$ of the further matching
edge determines at most one incoming matching.
\end{lemma}
\begin{proof}
By \noderef{MatchingThreeEdgeSwitchProperties},
$g(b_1)=f_1(x)=y=f_2(x)=g(b_2)$. Injectivity of $g$ gives
$b_1=b_2$; denote the common element by $b$.
The same value identities recover the original functions at the
three modified elements:
\[
 f_j(x)=y,\qquad f_j(a)=g(x),\qquad f_j(b)=g(a).
\]
At every other $z$, that node says $f_j(z)=g(z)$.
These four cases exhaust $A$ and show that $f_1$ and $f_2$ have
identical values everywhere. Hence they are equal.
In the matching interpretation, the edges of $g$ on the cycle are
$xg(x)$, $ag(a)$, and $bg(b)=by$. The first and last meet the ends
of $xy$, while the edge at $a$ is the further edge. This verifies
the reconstruction description.
\end{proof}
